\section{Introduction}

Learning a task that spans hundreds of milliseconds requires assigning credit to synapses whose
influence on behaviour is delayed and passes through recurrent dynamics. The exact solution,
back-propagation through time \citep[BPTT;][]{werbos1990bptt}, needs information that a synapse
does not have: the future error and its path through every other neuron. Local alternatives keep,
at each synapse, an \emph{eligibility trace} of how the synapse could have changed its
postsynaptic cell. They combine it with a broadcast error or reward signal, in line with
three-factor plasticity \citep{gerstner2018}. Examples are RFLO \citep{murray2019rflo} and e-prop
\citep{bellec2020eprop} \citep[see also][]{williams1989rtrl,marschall2020}. These rules are exact
for a cell's own dynamics, and they approximate the rest of the network, for example with random
feedback \citep{lillicrap2016fa}.

What the trace must contain depends on the cell's own dynamics, and real cortical neurons are not
the leaky units used in most such work. Their responses depend on several timescales, and most
show spike-frequency adaptation \citep{benda2003}. For spiking neurons with a single adaptive
threshold, adaptation is known to support memory over behavioural timescales
\citep{bellec2018lsnn,salaj2021sfa}, and e-prop includes the adaptation variable in its
eligibility trace \citep{bellec2020eprop}. Whether this matters for neurons whose intrinsic
dynamics were \emph{fitted to recordings}, with many heterogeneous filters per cell, has been hard
to test: data-fitted single-cell models were usually not differentiable \citep{teeter2018glif}, and
large data-constrained network models have been trained with BPTT
\citep{billeh2020,chen2022v1}. A recent bank of differentiable generalized firing-rate (GFR)
models, fitted to more than a thousand patch-clamp recordings \citep{yeung2025gfr}, now makes the
test possible.

We built recurrent networks entirely from these fitted cells, keeping each cell's intrinsic
parameters frozen and making only the synaptic weights plastic. We asked three questions:
(i) how close a cell-local, forward-in-time rule comes to the exact gradient on working-memory
tasks, compared with a frozen recurrent reservoir; (ii) whether the rule must include each cell's
fitted spike-history filter in its eligibility trace; and (iii) whether the specific identity of
the cells (the natural pairing of each cell's parameters and their heterogeneity) or the type-based
wiring matters for learning, once confounds are removed.

The third question carries a trap that we found in our own earlier analyses. Shuffling parameters
across cells moves each cell away from the operating point it was calibrated to. In our networks
this roughly doubled firing rates and multiplied saturation, so an apparent effect of
``destroying biological identity'' could not be separated from a change of operating point. We therefore
recalibrate every manipulated bank with the same label-free procedure before comparing. We also
split cells by donor animal, compare learners at learning rates tuned on a separate development
cohort, treat the network as the unit of inference \citep{aarts2014}, fit every population-analysis
axis on trials independent of those it is evaluated on \citep{kriegeskorte2009}, and test all
hypotheses in a protocol locked before the held-out confirmation cohort was touched.

Four things came out of the held-out confirmation (1,680 networks).
\begin{itemize}
\item The local rule needs each cell's fitted spike-history term. Without it, accuracy at unseen
delays dropped by 7--17\pp, to about the level of a frozen reservoir. With it, the local
estimate stayed aligned with the exact gradient, and the gain grew with each cell's adaptation
strength.
\item The local rule still trails BPTT on tasks that combine memory with a later cue. The gap
comes from ignoring paths through other cells, not from random feedback.
\item Cell identity and type-level wiring had no detectable effect once operating points were
matched. Heterogeneity helped BPTT strongly but not the local rule.
\item The operating point shaped local learning in a task-dependent way.
\end{itemize}
